\begin{titlepage}
\centering
\vspace*{1.2in}
{\Huge\bfseries\color{knobink} The Universe\\[6pt] of Knobs\par}
\vspace{18pt}
{\large\itshape How to fit anything, why you shouldn't,\\ and the one dial nobody gets to touch\par}
\vspace{1.1in}
{\large C. P. Zimmerman\par}
\vfill
{\small First edition, November 2026\par}
{\small DOI 10.5281/zenodo.23990956\par}
\end{titlepage}

\thispagestyle{empty}
\vspace*{\fill}
{\small
\noindent\textbf{The Universe of Knobs}\\
First edition, 5 November 2026. Published on Zenodo as version 1.0, DOI 10.5281/zenodo.23990956, under a CC BY 4.0 licence.

\medskip\noindent
This book is a popular and technical companion to two working notes from an AI-assisted research programme: \emph{ACDM flavored fried chicken version of Galaxy Dynamics} (version 2.3, DOI 10.5281/zenodo.23962973; version 2.0, DOI 10.5281/zenodo.23958263) and \emph{A Galaxy Law Tied to Dark Energy, and a Map of What It Leaves Unexplained} (PAPER67; version 1.0, DOI 10.5281/zenodo.23883741; version 1.1, DOI 10.5281/zenodo.23966443). Every number quoted from those notes can be checked against the scripts and frozen criteria committed in the programme's public repository. Neither note is peer reviewed, and no outside group has reproduced their numbers.

\medskip\noindent
\textbf{Read this before anything else.} The coefficient $\kappa=\tfrac12$ at the centre of this book is \emph{fitted}, not derived. The galaxy law is Milgrom's modified dynamics in its simple form, and the idea of tying its scale to the dark energy density was proposed by Limbach, Psaltis and \"Ozel in 2008. The theory is not closed. It fails in galaxy clusters, it supplies nothing for the cosmic microwave background, and it fails several other tests listed in Chapter~12. Nothing in this book says the data favour this framework over \LCDM\ or over anything else. The fried chicken recipes were tuned to imitate the law; the author does not believe them, and a good imitation is evidence of freedom, not of truth.

\medskip\noindent
This is a picture book on purpose: 130 figures, drawn by \texttt{figures.py}, \texttt{figures\_math\_a.py}, \texttt{figures\_math\_b.py}, \texttt{figures\_picture\_a.py} and \texttt{figures\_picture\_b.py}, with worked examples computed by \texttt{worked\_numbers.py}, all in the book's folder. Illustrations are drawings, not data; cartoons and mock data are labelled on the figure. Jokes are the author's own and carry no physics.

\medskip\noindent
Typeset in Latin Modern with \LaTeX.
}
\vspace*{\fill}
\clearpage
